\section{深度学习基础回顾}

本节课由 Justin Johnson 主讲，他是 CS231N 的创始人之一，曾在 Stanford 攻读博士学位（2012--2018），之后在 Facebook AI Research 和 University of Michigan 从事深度学习与计算机视觉研究。本节课是课程前半段（深度学习基础）与后半段（视觉世界感知）的分水岭，我们首先回顾前几节课的核心内容，然后深入介绍卷积神经网络。

\subsection{图像分类与线性分类器}

\begin{figure}[H]
\centering
\includegraphics[width=0.95\textwidth]{slides-images/slide-002.jpg}
\caption{课程前半段回顾：从图像分类到优化的完整流程\protect\footnotemark}
\end{figure}
\footnotetext{Slides 第3页。}

图像分类是计算机视觉最核心的任务之一：给定一张输入图像（三维张量，$H \times W \times 3$），模型需要输出一组分数，表示该图像属于每个预定义类别的可能性。

线性分类器是最简单的模型形式：

$$f(x) = Wx + b$$

其中 $x$ 是展平后的像素向量（如 $32 \times 32 \times 3 = 3072$ 维），$W$ 是权重矩阵，$b$ 是偏置。每一行 $W_i$ 可以被视为一个\textbf{类别模板}——模型通过计算输入图像与每个模板的内积来得到分类分数。

\begin{knowledgebox}{线性分类器的两种视角}
\begin{itemize}
    \item \textbf{模板视角}：权重矩阵的每一行是一个与整张图像大小相同的模板。分类分数就是输入与模板的相似度。缺陷：每个类别只能有\textbf{一个模板}，无法处理同一类物体的多种外观（如红色汽车 vs 蓝色汽车）。
    \item \textbf{几何视角}：线性分类器在高维空间中用超平面划分不同类别的区域。缺陷：如果数据不是线性可分的，线性分类器无能为力。
\end{itemize}
\end{knowledgebox}

\subsection{损失函数与优化}

有了模型之后，我们需要\textbf{损失函数}来衡量当前参数的好坏。常用的分类损失包括 Softmax 交叉熵损失和 SVM Hinge 损失。

\begin{figure}[H]
\centering
\includegraphics[width=0.95\textwidth]{slides-images/slide-004.jpg}
\caption{优化景观：横轴代表所有可能的权重配置，纵轴为损失函数值，目标是找到谷底\protect\footnotemark}
\end{figure}
\footnotetext{Slides 第5页。}

优化算法负责在参数空间中搜索低损失的区域。常用算法包括随机梯度下降（SGD）、带动量的 SGD、RMSProp 和 Adam 等。

\begin{importantbox}{Adam 获得 ICLR 2025 时间检验奖}
就在本节课授课前一天，ICLR 2025 将``时间检验奖''（Test of Time Award）颁给了 Adam 优化器的论文。Adam 最初于 2015 年在 ICLR 发表，至今仍是深度学习中使用最广泛的优化器之一。
\end{importantbox}

\subsection{从线性到非线性：神经网络}

线性分类器的表达能力严重不足。解决方案：堆叠多个线性变换并在它们之间插入\textbf{非线性激活函数}。

$$f(x) = W_2 \cdot \sigma(W_1 x + b_1) + b_2$$

这就是最简单的两层神经网络。虽然代数变化很小（多了一个矩阵和一个非线性函数），但分类能力得到了质的飞跃。

\subsection{计算图与反向传播}

\begin{figure}[H]
\centering
\includegraphics[width=0.95\textwidth]{slides-images/slide-008.jpg}
\caption{计算图与反向传播：数据从左到右前向传播计算输出，梯度从右到左反向传播\protect\footnotemark}
\end{figure}
\footnotetext{Slides 第9页。}

\begin{importantbox}{反向传播的核心思想}
反向传播将\textbf{全局问题}（计算损失对所有参数的梯度）转化为\textbf{局部问题}——每个计算节点只需要：
\begin{enumerate}
    \item 前向：根据输入计算输出
    \item 反向：接收上游梯度，计算对自身输入的梯度（链式法则）
\end{enumerate}
节点不需要知道自己处于什么样的网络中、要解决什么问题。这种局部 API 设计使得我们可以像搭积木一样构建任意复杂的神经网络。
\end{importantbox}

\begin{warningbox}{梯度形状法则}
对于标量损失 $L$：
\begin{itemize}
    \item 上游梯度 $\frac{\partial L}{\partial \text{output}}$ 的形状 = output 的形状
    \item 下游梯度 $\frac{\partial L}{\partial \text{input}}$ 的形状 = input 的形状
\end{itemize}
这是因为 $L$ 是标量，梯度本质上是``如果微扰张量中某个元素，损失会怎么变化''。
\end{warningbox}

\subsection{深度学习通用配方}

这几节课建立了解决几乎所有深度学习问题的通用配方：

\begin{enumerate}
    \item 将问题编码为\textbf{张量输入/输出}
    \item 设计\textbf{计算图}（神经网络架构）
    \item 选择\textbf{损失函数}
    \item 用\textbf{梯度下降 + 反向传播}优化
\end{enumerate}

无论是图像分类、图像生成还是大语言模型，底层都遵循这个配方。

\subsection{本章小结}

前四节课覆盖了深度学习的所有基础组件：线性分类器定义了问题形式，损失函数衡量了解的质量，优化算法搜索参数空间，神经网络提供了强大的函数族，计算图和反向传播让梯度计算自动化。接下来我们将进入课程第二阶段，学习如何针对图像设计更好的网络架构。

%% ========== Part 2: 从特征工程到端到端学习 ==========

